% ===========================================================================
% princeton.cls -- the example
% ===========================================================================
%
% This file does two jobs.
%
%   1. It documents every class option, in the block below. Each option is
%      written out explicitly with its alternatives in the comment, so you
%      change one word to see a different configuration.
%
%   2. It exercises every public command and environment the class provides.
%      Compiling it and looking at the result tells you whether the class is
%      working.
%
% Build it with:
%
%     make                 both engines, output in build/
%     lualatex example     from this directory
%
% Options are mutually exclusive by nature: a document cannot be columns=one
% and columns=two at once. To check a configuration, edit the line and rebuild.
% The ones most worth flipping are theme, columns and titleformat.
%
% ===========================================================================

\documentclass[
  11pt,                   % passed through to the base class
  class=article,          % article | report | book -- the base class
  theme=classic,          % classic | engineering -- which visual identity
  titleformat=simple,     % simple | cv | fancy   -- heading style
  columns=one,            % one | two             -- two-column body from p2
  logo=running,           % default | first | running | none
  logocolor=color,        % color | black         -- full colour or one colour
  signature=engineering,  % engineering | university | none  (theme=engineering)
  fontset=auto,           % auto | princeton | free
  marginpos=left,         % left | right          -- which side the margin sits
  itemcolor=pu-gray,      % any defined colour    -- the itemize bullet
  itemsep=1ex,            % any length            -- spacing between items
  colorlinks=true,        % true | false          -- coloured links
  date=true,              % true | false          -- date in the footer
  titledate=false,        % true | false          -- \date under the title
  draft=false,            % true | false          -- DRAFT marker in the footer
]{princeton}

\usepackage{blindtext}    % filler text, for this example only

% --- Metadata -------------------------------------------------------------
% \author and \email feed the margin block that \makeinfo prints. \version
% adds a version stamp to the footer and to that block.
\author{Jürgen Hackl}
\email{hackl@princeton.edu}
\address{
  Princeton University\\
  59 Olden Street\\
  Princeton, NJ 08544
}
\version{example}

\pretitle{A Princeton}
\title{Class Example}
\subtitle{Every option and every command}
\titlesize{\LARGE}

\begin{document}

\maketitle

% The author block. Lives in the margin, or becomes a tinted aside once the
% document goes two-column.
\makeinfo

\tableofcontents

% Everything up to the first \section is front matter. With columns=two the
% switch to two columns fires at the first \section, so this is the last text
% set against the wide margin.
\subsection*{Front matter}

With \texttt{columns=two} the body switches to two columns at the first
\verb|\section|, so everything above this point keeps the wide margin. Write
\verb|\startcolumns| yourself if you want the seam somewhere else.

\begin{margindescription}
\item[Layout] one column here
\item[Margin] real, not an aside
\end{margindescription}

% A margin figure and a margin note at the same height would overprint:
% marginnote and marginfigure do not know about each other. Keep them apart.
\blindtext

\begin{marginfigure}
  \caption{A \texttt{marginfigure}. In two columns it becomes an ordinary
    column-width float}
  \includegraphics[width=\linewidth]{example-image-a}
\end{marginfigure}

\section{The margin apparatus}

Everything in this section lives in the margin when the document is set in one
column. In two columns there is no margin, so each one falls back to something
sensible without being renamed.

\subsection{margindescription}

\begin{margindescription}
\item[Created] 2026-09-05
\item[Purpose] exercising the class
\item[Renamed] nothing
\end{margindescription}

\blindtext

\subsection{customdescription}

The same list without the margin wrapper, for use inline.

\begin{customdescription}
\item[Field] a description list with orange labels
\item[Second] a second field
\end{customdescription}

\section{Figures and tables}

\subsection{Side captions}

\begin{figure}[ht]
  \sidecaption{A \texttt{figure} with \texttt{\textbackslash sidecaption}, which
    puts the caption in the margin}
  \includegraphics[width=\linewidth]{example-image-a}
\end{figure}

\blindtext

\begin{figure}[ht]
  \SideCaption{A \texttt{figure} with \texttt{\textbackslash SideCaption}, the
    class wrapper around it}
  \includegraphics[width=\linewidth]{example-image-b}
\end{figure}

\blindtext

\subsection{princetonfigure, both modes}

This environment exists for Org mode export. Mode \texttt{f} is an ordinary
figure, mode \texttt{m} a margin figure.

\begin{princetonfigure}[ht][f]
  \includegraphics[width=\linewidth]{example-image-a}
  \caption{\texttt{princetonfigure} in mode f}
\end{princetonfigure}

\blindtext

\begin{princetonfigure}[ht][m]
  \includegraphics[width=\linewidth]{example-image-a}
  \caption{\texttt{princetonfigure} in mode m}
\end{princetonfigure}

\blindtext

\subsection{Full-width floats}

\texttt{largefigure} and \texttt{largetable} widen into the margin in one
column, and span both columns in two. From Org mode they are reached with
\texttt{:float largefigure} or \texttt{:float largetable} in the block's
\texttt{ATTR\_LATEX} line.

\begin{largefigure}
  \caption{A \texttt{largefigure}}
  \includegraphics[width=\linewidth]{example-image-a}
\end{largefigure}

\begin{largetable}
  \caption{A \texttt{largetable}, and a summary of what the margin commands do
    in each layout}
  \centering
  \begin{tabular}{p{34mm}p{40mm}p{40mm}}
    \hline
    Command & One column & Two column \\
    \hline
    \verb|\makeinfo|        & margin block   & tinted aside \\
    \verb|margindescr.|     & margin list    & tinted aside \\
    \verb|marginfigure|     & margin float   & column float \\
    \verb|\sidecaption|     & margin caption & ordinary caption \\
    \verb|largetable|       & widened        & spans both columns \\
    \hline
  \end{tabular}
\end{largetable}

\blindtext

\section{Type, lists and maths}

\subsection{Lists}

\begin{itemize}
\item An itemize item. Its bullet takes the \texttt{itemcolor} option
\item A second item, spaced by \texttt{itemsep}
\end{itemize}

\begin{enumerate}
\item An enumerate item
\item A second one
\end{enumerate}

\subsection{Figure styles}

The identity module exposes the alternate figure sets of whichever face is in
use.

\begin{itemize}
\item Lining figures 0123456789
\item Oldstyle, \verb|\textos|: \textos{0123456789}
\item Tabular, \verb|\texttab|: \texttab{0123456789}
\item Display face, \verb|\princetondisplay|: {\princetondisplay Princeton}
\end{itemize}

\subsection{Mathematics}

Maths is matched to the text face, so it changes with the theme.

\[
  \int_0^{\infty} e^{-x^2}\,dx = \frac{\sqrt{\pi}}{2},
  \qquad \alpha\beta\gamma \in \mathbb{R},
  \qquad \mathbf{A}\mathbf{x} = \mathbf{b}.
\]

\subsection{Links}

A hyperlink: \href{https://engineering.princeton.edu}{Princeton Engineering}.
A URL: \url{https://www.princeton.edu}.

\section{CV commands}

These work under any heading style. Switch the heading option to CV in the
block at the top to see those headings too, where subsection titles move into
the margin.

% \cvmargin and \cvitem must BEGIN a paragraph. \cvlabel hangs its key left of
% the current position, so at the start of a line it lands in the gap between
% the margin and the text block, which is the intent. After a run-in heading it
% would hang from wherever that heading ended, in the middle of the sentence.
\subsection*{Positions}

\cvmargin{Jan 2023 ongoing}\cvinfo{Princeton} An entry with the date in the
margin and the institution flushed right.

\cvmargin{2020 to 2023}\cvinfoit{Somewhere else} A second entry, with the
institution in italics.

\subsection*{Keyed entries}

\cvitem{2025}{J33} An entry keyed with \verb|\cvitem|. The date goes in the
margin and the key hangs in the gap to the left of this text, where it reads as
a reference tag rather than as part of the sentence.

\cvitem{2024}{J32} A second keyed entry, so the two keys line up.

\subsection*{Referring to them}

\verb|\cvref| takes a bracket style and a size. The usual form is

\verb|\cvref[both]{J33}[\scriptsize]|

which gives \cvref[both]{J33}[\scriptsize]. The other styles bracket a run as a
whole: \verb|left| gives \cvref[left]{J33}[\scriptsize] and \verb|right| gives
\cvref[right]{J32}[\scriptsize]. With no style it is bare: \cvref{J33}.

\section{Full-bleed background}

\verb|\fullfigure| stamps an image across the whole page, behind the text, on
the current page only. It takes \texttt{height}, \texttt{align} and
\texttt{angle}.

\fullfigure[height=2in,align=b]{example-image-b}

\blindtext

\end{document}

%%% Local Variables:
%%% mode: latex
%%% TeX-master: t
%%% End:
